\section{Absolute volumes at every cusp count}\label{sec:absolute}

We now assemble the preceding estimates to prove Theorem~\ref{thm:intro-absolute}. The extraction theorem, Theorem~\ref{thm:sharp-extraction}, gives the explicit leading term with relative error $O(g^{-1})$, and also the quantitative volume ratios. The expansion of Theorem~\ref{thm:canonical-expansion} then gives all higher corrections and their uniform remainder. We treat positive cusp counts first and recover $n=0$ from the exact dilaton relation.

For $g\ge2$ and $n\ge1$, set $h=2g-1$, $k=n-1$, and choose $x=X(Y)$
by $h\Psi(x)=k$. At $k=0$ take $x=Y=0$. Define
\begin{equation}
 \begin{split}
 \widehat V_{g,n}&=
 (2\pi^2)^{3g-3+n}\frac{\mathcal J_{g,1}(Y)}{\pi}
 \left(\frac{k}{\mathrm{e} x}\right)^k
 \sqrt{\frac{k}{h\mathscr D\Psi(x)}}\quad(n\ge1),\\
 \widehat V_{g,0}&=
 \frac{\Gamma(2g-2)(4\pi^2)^{2g-3}}{\sqrt{\pi g}}.
 \end{split}
 \label{eq:AV1}
\end{equation}
At $k=0$, both $k$-dependent factors on the first line are defined by
their limits, which are one. All functions in this formula were defined
in \eqref{eq:TL32}, \eqref{eq:PF1}, and \eqref{eq:PF3}.
Since $r=2\Delta$, $Aar=I/\sqrt\Delta$ by \eqref{eq:TL32}, and
$\Theta=-1/\bar K_Y'(1)=-a/(r\,\partial_zK_Y(r))$ by \eqref{eq:PF1}, one has
$r^{1/2}\sqrt{A\Theta}=\varpi(Y)$ and $\mathcal Z=\sqrt2\,I$, with
$\varpi$ as in the introduction. Hence \eqref{eq:AV1} is the formula
\eqref{eq:intro-leading}, with $\mathscr D\Psi(x)=x\Psi'(x)$.

\begin{theorem}[The leading term]\label{thm:absolute}
There is a constant $C>0$ such that, for all sufficiently large $g$,
\begin{equation}
 \sup_{n\ge0}\left|\frac{V_{g,n}}{\widehat V_{g,n}}-1\right|
       \le\frac Cg.
 \label{eq:AV2}
\end{equation}
The order $g^{-1}$ is sharp for the leading model \eqref{eq:AV1}:
\begin{equation}
 \liminf_{g\to\infty}g\sup_{n\ge0}
       \left|\frac{V_{g,n}}{\widehat V_{g,n}}-1\right|
       \ge\frac7{12}-\frac{17}{6\pi^2}>0.
 \label{eq:AV3}
\end{equation}
\end{theorem}
\begin{proof}
For $n\ge1$, apply Theorem~\ref{thm:sharp-extraction} with $p=1$.
The coefficient identity \eqref{eq:TL50} gives
\begin{equation}
 V_{g,n}=(2\pi^2)^{3g-3+n}(n-1)!x^{-k}
     \frac{\mathcal J_{g,1}(Y)}\pi\,
       \mathfrak p_k\sqrt{\frac{k}{h\mathscr D\Psi(x)}}
       (1+O(g^{-1})),
 \label{eq:AV4}
\end{equation}
uniformly, including $k=0$ by continuity.
Since $k!\mathfrak p_k=(k/\mathrm{e})^k$, this is precisely the first line of
\eqref{eq:AV1} with uniform relative error $O(g^{-1})$.
For $n=0$, the constant was independently identified in
Section~\ref{sec:prefactor}. The fixed-cusp expansion of
Mirzakhani--Zograf~\cite[Theorem~1.2 and Remark~1.3]{MZ} yields
\begin{equation}
 \frac{V_{g,0}}{\widehat V_{g,0}}
 =1+\frac1g\left(\frac7{12}-\frac{17}{6\pi^2}\right)+O(g^{-2}).
 \label{eq:AV5}
\end{equation}
This closes \eqref{eq:AV2} and proves \eqref{eq:AV3}.
The latter concerns the uncorrected model chosen here, rather than an
obstruction to adding further asymptotic corrections.
\end{proof}

The model \eqref{eq:AV1} is explicit and contains no unknown constant.
Its action $I(Y)$ and the spectral-edge parameter $Y$ have the
matrix-model interpretation recalled in Subsection~\ref{subsec:jt-edge};
what is proved here is the uniform relative comparison with the true
volumes. The normalization is fixed at the pure-$\psi$ endpoint
(Section~\ref{sec:prefactor}); the fixed-cusp expansion of~\cite{MZ}
enters only afterwards, for $n=0$ in the proof above.

\subsection{An inverse-genus error for the volume ratios}
\label{subsec:quantitative-volume-ratios}

The sharper extraction estimate also gives a uniform rate in the
volume-ratio theorem. The two genus-exchange generating functions have
the same saddle. Adding a cusp changes the complexity by one; the
following argument controls the resulting saddle displacement even
near the endpoints.

\begin{proof}[Completion of the proof of Theorem~\ref{thm:volume}]
For $n\ge1$ put $h=2g-1$, $k=n-1$, and $h\Psi(x)=k$.
Apply Theorem~\ref{thm:sharp-extraction} at $(g,p)=(g,1)$ and
$(g-1,3)$. Their Poisson and variance factors cancel exactly, and
\eqref{eq:PF3} gives $\mathcal J_{g-1,3}/\mathcal J_{g,1}=4r$.
Consequently \eqref{eq:TL50} yields
\[
 B_{g,n}=\mathcal B(2k/h)(1+O(h^{-1})).
\]

For the cusp ratio assume first that $k\ge1$, and let $x'$ satisfy
$(h+1)\Psi(x')=k$. Write $q=\log x$, $q'=\log x'$ and
$v=\mathscr D\Psi$. The convergent endpoint expansions give
\[
 \sup_{0<x<x_c}\left\{
 \frac{\Psi(x)^2}{v(x)}+
 \left|\frac{\Psi(x)}{v(x)}\mathscr D\log v(x)\right|
 +|\beta_2(x)|\right\}<\infty.
\]
Indeed, near zero $\Psi\sim v\sim2x$; near criticality
$\Psi\sim5/(4\varepsilon)$ and $v\sim5/(4\varepsilon^2)$.
The function $\beta_2$ is bounded by \eqref{eq:PF7}.
Since $\Psi(x')/\Psi(x)=h/(h+1)$, integration with respect to
$\log\Psi$ shows that
\[
 \frac{v(x')}{v(x)}=1+O(h^{-1}),\qquad
 |q'-q|\le\frac{C\Psi(x)}{h v(x)}.
\]
For $L(q)=\log\mathcal J_{g,2}(\mathrm{e}^q)-kq$, the exact model identities give
\[
 L'(q)=(h+1)\Psi(\mathrm{e}^q)-k+\beta_2(\mathrm{e}^q)\Psi(\mathrm{e}^q).
\]
On $[q',q]$ its absolute value is at most $C\Psi(x)$, so
$|L(q')-L(q)|\le C\Psi(x)^2/(hv(x))=O(h^{-1})$.
Thus the saddle displacement and the variance factor in
\eqref{eq:CE2} each change the coefficient ratio by $1+O(h^{-1})$.
At the common point $x$, \eqref{eq:PF3} gives exactly
\[
 \frac{\mathcal J_{g,2}(x)}{\mathcal J_{g,1}(x)}
 =\frac{h}{Aar}\sqrt{\frac{h}{h+1}}.
\]
Using the coefficient normalization in \eqref{eq:TL53}, we obtain
\[
 C_{g,n}=\frac{h}{2g}\mathcal C(2k/h)(1+O(h^{-1})).
\]
At $k=0$ both saddles are zero. The conventions in
Theorem~\ref{thm:sharp-extraction} give the same two ratio formulas
directly, without logarithmic coordinates.

The positive profiles are analytic at zero, and their expansions at
infinity in \eqref{eq:SP13}--\eqref{eq:SP15} are convergent in the
inverse square root of the density. Hence
\[
 \sup_{\lambda\ge0}(1+\lambda)
 \bigl(|(\log\mathcal B)'(\lambda)|+
       |(\log\mathcal C)'(\lambda)|\bigr)<\infty.
\]
For $\lambda_*=2(n-1)/(2g-1)$ and $\lambda=n/g$,
$|\lambda_* -\lambda|/(1+\min(\lambda_*,\lambda))=O(g^{-1})$
uniformly in $n\ge1$. Integrating the preceding bound replaces
$\lambda_*$ by $\lambda$ with relative error $O(g^{-1})$.
For $n=0$, the fixed-cusp expansion of~\cite{MZ} at the finitely
many cusp counts $0,1,2$ gives the required rate directly.
Finally $D_{g,n}=C_{g,n-1}/C_{g,n}$ and the same logarithmic
derivative bound prove $\sup_{n\ge1}|D_{g,n}-1|=O(g^{-1})$.
\end{proof}

\subsection{Proof of Theorem~\ref{thm:intro-absolute}}

For $n\ge1$ the expansion is obtained by applying the extraction
theorem, Theorem~\ref{thm:canonical-expansion}, with $p=1$, keeping the
same leading term. The unmarked case $n=0$ is recovered from $n=1$ by the
exact dilaton relation, so its higher coefficients need no separate
asymptotic normalization. For $h=2g-1$, the scalar $\mathsf C_0$ of \eqref{eq:FE1} obeys
\begin{equation}
 V_{g,0}=\frac{V_{g,1}\mathsf C_0(h-1;0)}{4\pi^2(h-1)},\qquad
 \frac{\widehat V_{g,1}}{4\pi^2(h-1)\widehat V_{g,0}}
       =\sqrt{1+\frac1h}.
 \label{eq:AV6}
\end{equation}
Indeed \eqref{eq:TL12} holds at $p=0$; at $Y=0$, undo its
normalizations with $r=\pi^2/2$, $A=1/\pi^2$.
The scalar expansion of $\mathsf C_0$ follows directly from the last
line of \eqref{eq:FE5}, using the proved one-marking profile and its
summable remainder. This step does not require an unmarked harmonic
comparison. The second identity follows from \eqref{eq:AV1} exactly.

The coefficients in Theorem~\ref{thm:intro-absolute} are defined as
follows. For $n=0$, define constants $a_{0,j}$ by the finite-order formal
rule
\begin{equation}
 1+\sum_{j\ge1}a_{0,j}t^j
 =(1+t)^{1/2}\mathsf C_0\!\left(\frac t{1-t};0\right)
       \left(1+\sum_{j\ge1}\mathcal L_{1,j}(0)t^j\right).
 \label{eq:AV7}
\end{equation}
Here the second argument denotes $Y=0$, and $t=1/h$; thus the
cusp scalar uses complexity $h-1$. Put $A_j(g,0)=a_{0,j}$, and, for
$n\ge1$, put $A_j(g,n)=\mathcal A_{1,j}(Y_{g,n})$ with the coefficients
\eqref{eq:CG3} of Theorem~\ref{thm:canonical-expansion}, where
$(2g-1)\Psi(X(Y_{g,n}))=n-1$.

\begin{proof}[Proof of Theorem~\ref{thm:intro-absolute}]
For $n\ge1$, take $p=1,k=n-1$ in
Theorem~\ref{thm:canonical-expansion}. The coefficient identity \eqref{eq:TL50}
and $k!\mathfrak p_k=(k/\mathrm{e})^k$ identify $[z^k]\Phi_g^{(1)}/\mathcal Q_{g,1,k}$
with $V_{g,n}/\widehat V_{g,n}$, including $k=0$.
For $n=0$, the dilaton scalar in \eqref{eq:AV6} has a one-marking
child of complexity $h$, so exactly
\[
 \mathsf C_0(h-1;0)=\frac1{A(0)}
       \sum_{d\ge1}b_d(0)c_{g,1}(d;0).
\]
Apply \eqref{eq:FE7} at this child with $L=K+1$.
Choose its fixed $M$ large enough for the absolute exponential
$b_d$ moment. With $S=1+ML$,
\[
 \sum_d\frac{|b_d(0)|}{A(0)}(S+d)^{2L}
 \le S^{2L}\sum_d\frac{|b_d(0)|}{A(0)}\mathrm{e}^{2d/M}
 \le C D^L(L!)^2.
\]
This gives a Gevrey-2 remainder for the whole scalar sum, including
invalid degrees. Its coefficient series is precisely
$\mathsf C_0(t/(1-t);0)$: the child substitution in \eqref{eq:FE5}
becomes $(t/(1-t))/(1+t/(1-t))=t$.

The other two exact factors in \eqref{eq:AV6} are
$\sqrt{1+1/h}$, with a convergent expansion, and the $n=1$ ratio,
with the marked estimate at zero. All three exact factors are bounded.
The finite-product identity used in Proposition~\ref{prop:hierarchy}
contains one remainder per term; its factorial weights obey
$(j!)^2((L-j)!)^2\le(L!)^2$. Applied twice, it gives the same
bound for their product and exactly the coefficients in \eqref{eq:AV7}.
Take common constants with the $n\ge1$ case. The coefficient bound
$|A_j(g,n)|\le CB^j(j!)^2$ follows from those of
Theorem~\ref{thm:canonical-expansion} and Proposition~\ref{prop:hierarchy}.
Finally, for $K=\lfloor c'\sqrt{2g-1}\rfloor$ with $c'\le c$ chosen after
the final base $B$ so that $B(c')^2<1$, the inequality
$B^L(L!)^2h^{-L}\le(BL^2/h)^L$ gives the exponentially small error.
\end{proof}

In particular, the explicit first correction \eqref{eq:CG12}
gives a relative error $O(g^{-2})$ uniformly over all cusp counts.
At the unmarked endpoint $a_{0,1}=7/6-17/(3\pi^2)$, consistently
with \eqref{eq:AV5} after $h=2g-1$.
The fixed-cusp existence of an all-order expansion is the theorem
of Mirzakhani--Zograf~\cite{MZ}; here the coefficients and remainders
remain controlled as the cusp count crosses every density scale.
